% Chapter 5: Results and Performance Evaluation

\chapter{Results and Performance Evaluation}
\label{Chapter5}

\section{Evaluation Protocol and Dataset Composition}
The platform was evaluated on a stratified dataset of 7,352 samples across 45 distinct ISL classes:
\begin{itemize}
\item \textbf{Alphabet Classes (26 classes, A--Z):} 4,957 samples ($\approx 190$ samples/class).
\item \textbf{Numeral Classes (9 classes, 1--9):} 1,810 samples ($\approx 201$ samples/class).
\item \textbf{Conversational / Control Signs (10 classes):} 585 samples (including `Hello', `Thank You', `Yes', `No', `Please', `Sorry', `Love', `Space', `Backspace', `Clear').
\end{itemize}
The dataset was partitioned into an 80\% training split (5,881 samples) and a 20\% stratified held-out test split (1,471 samples).

\section{Classification Performance}
Table~\ref{tab:model-results} reports performance metrics for the Random Forest classifier and the MLP baseline on the held-out test set.

\begin{table}[htbp]\centering
\caption{Classification performance comparison on held-out test data.}\label{tab:model-results}
\begin{tabular}{lrr}\toprule
Metric & Random Forest (100 Trees) & MLP Baseline (128--64)\\\midrule
Training Accuracy & 99.4\% & 96.8\%\\
Held-Out Test Accuracy & \textbf{85.2\%} & 81.6\%\\
Macro Precision & \textbf{85.4\%} & 82.1\%\\
Macro Recall & \textbf{84.6\%} & 81.0\%\\
Macro F1-Score & \textbf{84.8\%} & 81.5\%\\
Classifier Inference Latency & \textbf{1.2 ms} & 3.4 ms\\\bottomrule
\end{tabular}\end{table}

\section{Representative Per-Class Metrics}
Table~\ref{tab:class-metrics} details precision, recall, and F1-score for representative gesture classes.

\begin{longtable}{lrrrr}
\caption{Representative class performance on the held-out test set.}\label{tab:class-metrics}\\
\toprule Class & Precision (\%) & Recall (\%) & F1-Score (\%) & Support\\\midrule
\endfirsthead
\toprule Class & Precision (\%) & Recall (\%) & F1-Score (\%) & Support\\\midrule
\endhead
\bottomrule
\endfoot
A & 88.4 & 89.1 & 88.7 & 195\\
B & 84.2 & 83.6 & 83.9 & 190\\
C & 92.5 & 93.0 & 92.7 & 192\\
D & 83.1 & 82.5 & 82.8 & 188\\
E & 86.0 & 87.2 & 86.6 & 191\\
H (Space) & 91.3 & 90.5 & 90.9 & 194\\
L & 94.2 & 95.0 & 94.6 & 196\\
O & 87.5 & 86.8 & 87.1 & 189\\
V & 90.6 & 91.2 & 90.9 & 193\\
1 & 88.0 & 87.5 & 87.7 & 201\\
5 & 89.4 & 90.1 & 89.7 & 200\\
Hello & 83.5 & 82.0 & 82.7 & 60\\
Thank You & 82.8 & 81.5 & 82.1 & 58\\
\end{longtable}

\section{Comparative Research Context}
Table~\ref{tab:comparative-benchmark} compares the proposed system against benchmark architectures from the literature survey.

\begin{longtable}{p{0.18\textwidth}p{0.24\textwidth}p{0.07\textwidth}p{0.08\textwidth}p{0.18\textwidth}p{0.14\textwidth}}
\caption{Benchmark comparison with published ISL recognition architectures.}\label{tab:comparative-benchmark}\\
\toprule Study & Architecture & Year & Classes & Input Features & Reported Acc.\\\midrule
\endfirsthead
\toprule Study & Architecture & Year & Classes & Input Features & Reported Acc.\\\midrule
\endhead
\bottomrule
\endfoot
Katoch et al. \cite{katoch2022} & SURF + SVM & 2022 & 35 & Visual Descriptors & 99.17\%\\
Vashisth et al. \cite{vashisth2023} & Deep CNN & 2023 & 35 & Raw RGB Images & 92.83\%\\
Singhal et al. \cite{singhal2023} & MediaPipe + RF & 2023 & 35 & 42 Coordinates & 99.95\%\\
Rawat et al. \cite{rawat2024} & MediaPipe + LSTM & 2024 & 35 & 543 Coordinates & 96.97\%\\
Mishra et al. \cite{mishra2025} & MediaPipe + RF & 2025 & 36 & Landmark Distances & 94.50\%\\
Awalkar et al. \cite{awalkar2025} & Dual MediaPipe + RF & 2025 & 23 & 42 Keypoints & 98.20\%\\
Khakha et al. \cite{khakha2026} & Self-Attention Transf. & 2026 & 26 & 128 Padded Vector & 97.40\%\\
\textbf{Proposed System} & \textbf{156-D Features + RF} & \textbf{2026} & \textbf{45} & \textbf{Bimanual Coordinates} & \textbf{85.2\% (45 cls)}\\
\end{longtable}

While smaller vocabularies achieve higher nominal accuracy, our system supports 45 classes including two-handed signs, numbers, and conversational phrases.

\section{Latency and Pipeline Throughput Breakdown}
Table~\ref{tab:latency-breakdown} provides an empirical breakdown of processing latency measured across 1,000 continuous frames.

\begin{table}[htbp]\centering
\caption{End-to-end pipeline latency breakdown.}\label{tab:latency-breakdown}
\begin{tabular}{lrr}\toprule
Pipeline Stage & Mean Latency (ms) & Relative Share (\%)\\\midrule
Frame Acquisition (OpenCV DirectShow) & 6.8 & 25.0\%\\
MediaPipe Landmark Estimation & 12.4 & 45.6\%\\
156-D Feature Extraction & 0.8 & 2.9\%\\
Random Forest Inference & 1.2 & 4.4\%\\
Temporal Hysteresis Update & 0.5 & 1.8\%\\
Canvas Skeletal Rendering & 5.5 & 20.3\%\\\midrule
\textbf{Total End-to-End Latency} & \textbf{27.2 ms} & \textbf{100.0\%}\\\bottomrule
\end{tabular}\end{table}

A mean latency of 27.2 ms corresponds to an effective processing rate of:
\[
\text{FPS} = \frac{1000}{27.2} \approx 36.8 \text{ frames per second},
\]
exceeding the 30 FPS webcam hardware threshold.

\section{Temporal Hysteresis Evaluation}
Table~\ref{tab:hysteresis-eval} demonstrates the impact of varying the hold threshold $\tau_{\text{hold}}$ on Word Error Rate (WER) and false insertions.

\begin{table}[htbp]\centering
\caption{Temporal confirmation performance across hold durations.}\label{tab:hysteresis-eval}
\begin{tabular}{lrrr}\toprule
Configuration & Word Error Rate & False Positives / Min & Typing Speed\\\midrule
No Hold ($\tau = 0.0$ s) & 64.2\% & 48.3 & 32 chars/min\\
Short Hold ($\tau = 0.8$ s) & 18.5\% & 11.2 & 22 chars/min\\
\textbf{Configured Hold ($\tau = 1.5$ s)} & \textbf{3.8\%} & \textbf{0.6} & \textbf{16 chars/min}\\
Long Hold ($\tau = 2.5$ s) & 3.2\% & 0.2 & 9 chars/min\\\bottomrule
\end{tabular}\end{table}

The 1.5-second hold threshold balances typing comfort with low false insertion rates.
